\chapter{Spacetime Rendered: General Relativity as a Geometric Codec}

In the previous chapter we proved that the entire particle spectrum of
the Standard Model organises into fermion complexity classes $(n, m)$,
with no gauge groups, coupling constants, or postulated force carriers
required. Bosons are compression residuals. Fermions are pixels. The
conservation law $\|F\|^2 + \|B\|^2 = 1$ holds exactly for any pure
state.

This chapter asks whether the same decomposition applies at the
opposite end of the scale hierarchy --- to spacetime geometry itself.
The answer is yes, and the formula is identical.

\section{The Same Operation, Different Degrees of Freedom}

The central operation of the codec framework is:

\[
\rho \;\longmapsto\;
\underbrace{\mathrm{diag}(\rho)}_{\text{local}}
+
\underbrace{\rho - \mathrm{diag}(\rho)}_{\text{non-local}}
\]

Applied to fermionic configurations, the diagonal gives Pauli
exclusion and the off-diagonal gives the photon, the gluon, and the
Born rule. This was established in the previous two chapters.

The present chapter applies the identical operation to metric
configurations --- pairs of conformal factors specifying how space is
stretched at each point on a discrete spatial grid. The diagonal part
will give local curvature, the source term of the Einstein equations.
The off-diagonal part will give tidal forces and gravitational waves.
Neither is postulated. Both follow from the definition
$B = \rho - \mathrm{diag}(\rho)$ and the geometry of three-dimensional
space.

One important scope note before proceeding. The theorems in this
chapter are proved within a discrete toy model: conformal metric
configurations $g(x) \in \mathbb{R}$ on a finite spatial grid of $L$
sites, encoded as $\psi = g^{(0)} + ig^{(1)}$. This captures the
qualitative structure of the Ricci/Weyl decomposition and the graviton
amplitude exactly. It does not yet recover the full rank-4 tensorial
Riemann curvature in four continuous dimensions. The extension to the
Einstein field equations in the continuum limit is the primary open
problem, stated at the end of this chapter.

\section{Theorem 1: The Ellipse as Minimum-Complexity Worldline}

Before decomposing metric configurations, we ask a more basic question:
what is the most compressible closed trajectory through spacetime?

A closed worldline $\psi(t)$ satisfying $\psi(0) = \psi(T)$ must
have a Fourier spectrum restricted to integer multiples of
$\omega_0 = 2\pi/T$ --- only these modes satisfy the periodicity
condition. The spectral complexity cost of mode $k$ is $|k|$ in units
of $\Delta\omega$. The trivial stationary solution has $C_s = 0$ but
represents no motion. The minimum nonzero cost is $|k| = 1$, achieved
by a single complex exponential:

\[
\psi(t) = A\,e^{i\omega_0 t}.
\]

In the complex plane, a single exponential traces an ellipse. Any
closed trajectory requiring additional Fourier modes costs $C_s \geq 2$.

Under Solomonoff induction, the ellipse carries probability weight
$2^{-1} = 1/2$. All other closed trajectories are exponentially
suppressed. Keplerian orbits are selected by minimum description
length, not by a force law.

The innermost stable circular orbit of a Schwarzschild black hole
falls out as a corollary: it is the minimum-complexity stable closed
worldline, and substituting $r_{\mathrm{ISCO}} = 6M$ into Kepler's
third law gives $\omega^2_{\mathrm{ISCO}} \cdot r^3_{\mathrm{ISCO}} = M$
exactly in Planck units. The ISCO is selected without postulating GR.

\section{Theorem 2: The Ricci/Weyl Decomposition}

A two-frame metric pair $(g^{(0)}, g^{(1)})$ is encoded as
$\psi(x) = g^{(0)}(x) + ig^{(1)}(x)$, normalised to unit norm.
Applying the standard density matrix decomposition:

\[
\rho = \rho_{\mathrm{Ricci}} + \rho_{\mathrm{tidal}} +
\rho_{\mathrm{graviton}},
\]

where

\[
\rho_{\mathrm{Ricci}} = \mathrm{diag}(\rho), \qquad
\rho_{\mathrm{tidal}} = \tfrac{1}{2}(\rho_W + \rho_W^\top), \qquad
\rho_{\mathrm{graviton}} = \tfrac{1}{2}(\rho_W - \rho_W^\top),
\]

and $\rho_W = \rho - \mathrm{diag}(\rho)$.

The three components are mutually orthogonal and their traces satisfy:
$\mathrm{Tr}(\rho_{\mathrm{Ricci}}) = 1$,
$\mathrm{Tr}(\rho_{\mathrm{tidal}}) = 0$,
$\mathrm{Tr}(\rho_{\mathrm{graviton}}) = 0$.

The physical identification maps precisely onto the Riemann tensor
decomposition of GR:

\begin{center}
\begin{tabular}{lll}
\toprule
Component & Riemann analogue & Physical role \\
\midrule
$\rho_{\mathrm{Ricci}}$ & Ricci tensor $R_{\mu\nu}$ &
Local curvature; matter source \\
$\rho_{\mathrm{tidal}}$ & Coulomb Weyl & Tidal forces; static \\
$\rho_{\mathrm{graviton}}$ & Radiative Weyl & Gravitational waves \\
\bottomrule
\end{tabular}
\end{center}

The tracelessness of both Weyl components is not imposed. It follows
from the trace-zero property of off-diagonal density matrices ---
the same mathematical fact that gives the tracelessness of the Weyl
tensor in GR, where it follows from the symmetries of the Riemann
tensor. Here it is a codec property.

The decomposition also gives two exact conditions. The graviton
component vanishes if and only if $g^{(0)} \propto g^{(1)}$ --- a
static metric configuration with no wave. The tidal component vanishes
if and only if $\psi$ is a single-site state with no spatial extent.
Both conditions are exact and encoding-independent.

\section{Theorem 3: The Graviton Amplitude}

Consider a conformal metric pair interpolating between two orthogonal
normalised metric profiles $\tilde{h} \perp \tilde{k}$:

\[
\psi = \cos\epsilon\,\tilde{h} + i\sin\epsilon\,\tilde{k},
\quad \epsilon \in [0, \pi/2].
\]

The graviton amplitude $\|W(\epsilon)\| = \|\rho_{\mathrm{graviton}}\|_F$
satisfies:

\[
\boxed{\|W(\epsilon)\| = \frac{\sin 2\epsilon}{\sqrt{2}}.}
\]

This is independent of the metric profiles $\tilde{h}$, $\tilde{k}$,
the number of sites $L$, and any overall scale.

The proof is identical in structure to the two-site fermion proof of
the previous chapter. The off-diagonal antisymmetric entries are
$-\tfrac{i}{2}\sin 2\epsilon\,[\tilde{h}(x)\tilde{k}(y) -
\tilde{k}(x)\tilde{h}(y)]$, and summing their squared magnitudes over
all pairs $(x, y)$ with $\langle\tilde{h}, \tilde{k}\rangle = 0$ gives
$\|W\|^2 = \sin^2(2\epsilon)/2$.

The graviton lifecycle traces the same arc as the virtual photon:
$\|W\| = 0$ at $\epsilon = 0$ (static metric), $\|W\| = 1/\sqrt{2}$
at $\epsilon = \pi/4$ (peak gravitational wave amplitude),
$\|W\| = 0$ at $\epsilon = \pi/2$ (static again). This is the
gravitational wave propagator, derived without postulating a graviton
field.

The formula $\sin(2\epsilon)/\sqrt{2}$ is not a coincidence shared
between the photon and graviton results. It is the unique consequence
of the density matrix decomposition applied to any normalised
configuration pair in superposition, regardless of whether those
configurations are fermionic or metric. The formula is universal
because the operation is universal.

\section{Two Polarisations Without a Spin-2 Postulate}

The two graviton polarisation states ($+$ and $\times$) emerge from
the eigenstructure of $\rho_{\mathrm{tidal}}$. For a wave metric pair
with alternating profiles, the eigenvalues of $\rho_{\mathrm{tidal}}$
include two degenerate values $\lambda = -1/4$, with eigenvectors:

\[
v_+ = \tfrac{1}{\sqrt{2}}[0, 1, 0, -1], \qquad
v_\times = \tfrac{1}{\sqrt{2}}[1, 0, -1, 0].
\]

These are the discrete analogues of alternating compression and
expansion along perpendicular axes --- the defining signature of
gravitational wave polarisation. The degeneracy is protected by the
reflection symmetry of the metric profiles, the lattice analogue of
the rotational symmetry that protects graviton polarisation in GR. No
spin-2 field is postulated.

\section{Theorem 4: The Newtonian Potential from Flux Conservation}

The graviton amplitude theorem gives $\|W\|_{\max} = 1/\sqrt{2}$ at
peak emission. In three spatial dimensions, the graviton propagates
outward from its source. By conservation of graviton flux over
spherical shells of area $4\pi r^2$:

\[
A(r) = \frac{\|W\|_{\max}}{2\sqrt{\pi}\,r} = \frac{1}{2\sqrt{2\pi}\,r}.
\]

The field amplitude falls exactly as $1/r$. This is exact and follows
from spherical geometry alone.

For two masses $M$ and $m$ separated by distance $R$, computing the
overlap integral of their graviton fields and assuming $\|W\| \propto M$
(more bits produce stronger field --- an assumption stated honestly as
Open Problem 1):

\[
V(R) = -\int_{\ell_P}^{\infty} A_M(r)\,A_m(r)\,4\pi r^2\,\mathrm{d}r
= -\frac{Mm}{8\pi R}.
\]

Matching to the Newtonian form $V(R) = -GMm/R$:

\[
\boxed{G = \frac{1}{8\pi} \quad \text{(Planck units).}}
\]

The factor $8\pi = 2 \times 4\pi$ arises entirely from spherical
geometry. It is the same $4\pi$ that appeared in Chapter~4 when
recovering the Bekenstein--Hawking entropy, and the same $4\pi$ that
converts flat raster area to spherical surface area. It is not a free
parameter. It is the unique conversion factor between discrete counting
on a flat grid and propagation through three-dimensional spherical
space, and it appears wherever the framework makes that transition.

The $1/r^2$ force law follows analytically from $V \propto 1/R$.
No force law was postulated. Gravity emerges from flux conservation.

\section{The Conservation Law}

The pure-state identity $\|\rho\|^2 = 1$ gives a metric-configuration
conservation law directly:

\[
\|\rho_{\mathrm{Ricci}}\|^2 + \|\rho_{\mathrm{Weyl}}\|^2 = 1,
\]

where $\|\rho_{\mathrm{Weyl}}\|^2 = \|\rho_{\mathrm{tidal}}\|^2 +
\|\rho_{\mathrm{graviton}}\|^2$.

This is the metric-configuration analogue of $\|F\|^2 + \|B\|^2 = 1$
from the previous chapter. Both are instances of the same Pythagorean
identity on Hilbert space. The bit budget spent on local curvature and
the bit budget spent on propagating waves must sum to one. You can
concentrate information into a dense mass, or you can let it radiate
outward as gravitational waves. The total is conserved.

Quantum mechanics and general relativity share not only the formula
$\sin(2\theta)/\sqrt{2}$ but the same conservation structure. They
are two projections of a single compression principle onto different
physical degrees of freedom.

\section{$G$ as a Consistency Condition}

Standard physics treats $G$ as a dimensionful coupling constant
measured by experiment. The framework gives a different account.

In Planck units, $G = 1$ by definition. Theorem~4 gives $G = 1/(8\pi)$
from the graviton flux argument. These are consistent: the
normalisation $\|W\| \propto M$ is not yet fixed independently of $G$.
What is fixed is the functional form --- $V \propto -Mm/R$ with the
$1/R$ dependence exact and the $8\pi$ geometric factor exact.

The value of $G$ in SI units is a unit conversion:

\[
G_{\mathrm{SI}} = \frac{\ell_P^3}{m_P\,t_P^2}.
\]

$G = 1$ in Planck units is the statement that the minimum-complexity
orbital geometry and the minimum-complexity orbital frequency are
mutually consistent at the Planck scale. $G$ is not a free parameter
of the universe. It is a consistency condition between the geometric
and temporal resolutions of the 184-bit system.

\section{What Remains Open}

Four open problems are worth stating precisely.

\textbf{Derivation of $\|W\| \propto M$.} Theorem~4 assumes the
graviton amplitude scales linearly with mass. This should follow from
the counting equation $n_{\mathrm{bh}} = \log_2(r_s/\ell_P) = 1 +
\log_2 M$, since more bits produce more off-diagonal weight. Making
this rigorous would complete Theorem~4.

\textbf{The continuum limit and the Einstein equations.} All theorems
here are proved on a discrete grid of conformal scalar factors. The
Einstein equations should emerge as the large-deviation stationarity
condition of $C_s$ over metric configurations in the continuum limit,
analogous to how the Euler--Lagrange equations are stationarity
conditions of the action.

\textbf{Full tensorial Riemann curvature.} The toy model uses scalar
conformal factors, not the full rank-4 Riemann tensor. Extending the
codec decomposition to the tensorial case in four dimensions is
required for a complete derivation of GR.

\textbf{Electromagnetic force law.} The same flux argument applied to
the $n = 2$, $m = 1$ fermion's antisymmetric residual should give the
Coulomb potential $V_{\mathrm{EM}}(r) \propto 1/r$. The ratio
$V_{\mathrm{EM}}/V_{\mathrm{grav}}$ would then address the hierarchy
problem from within the framework. This derivation is not yet complete.

\section{What This Chapter Establishes}

Four theorems, each exact within the conformal metric toy model.

The ellipse theorem: Keplerian orbits are selected by Solomonoff
induction. The minimum-complexity closed worldline is a single Fourier
mode, tracing an ellipse. No force law is required.

The curvature decomposition theorem: $\rho = \rho_{\mathrm{Ricci}} +
\rho_{\mathrm{tidal}} + \rho_{\mathrm{graviton}}$ exactly, with both
Weyl components traceless and two graviton polarisation modes emerging
as degenerate eigenmodes of $\rho_{\mathrm{tidal}}$.

The graviton amplitude theorem: $\|W(\epsilon)\| =
\sin(2\epsilon)/\sqrt{2}$, identical in form to the Born rule identity
of Chapter~7. The conservation law $\|\rho_{\mathrm{Ricci}}\|^2 +
\|\rho_{\mathrm{Weyl}}\|^2 = 1$ mirrors $\|F\|^2 + \|B\|^2 = 1$ of
Chapter~8.

The Newtonian potential theorem: graviton flux conservation over
$4\pi r^2$ shells gives $A(r) \propto 1/r$ exactly, and
$V(R) = -GMm/R$ with $G = 1/(8\pi)$ in Planck units. The factor
$8\pi$ is the same spherical geometry factor that appeared when
recovering Bekenstein--Hawking entropy in Chapter~4.

Quantum mechanics and general relativity are two projections of the
same compression principle. The local/non-local split of the density
matrix produces Pauli exclusion and photons when applied to fermions,
and Ricci curvature and gravitational waves when applied to metric
configurations. The formula, the conservation structure, and the
geometric factor are the same in both cases.

The derivation of the Einstein field equations in the continuum limit
is the primary outstanding goal of the research programme.
